\section{The test}\label{sec:the-test}

\begin{definition}[Roles]
The low individual degree test will be played between two provers and a verifier.
The two provers are named Player~$\mathrm{A}$ and Player~$\mathrm{B}$.
A \emph{role} $r$ is an element of the set $\{\mathrm{A}, \mathrm{B}\}$.
Given a role~$r$, we write $\overline{r}$ for the other element of the set $\{\mathrm{A}, \mathrm{B}\}$.
\end{definition}

\begin{definition}[Low individual degree test]
Let~$m$ and~$d$ be nonnegative integers.
Let $q$ be a prime power.
Then the \emph{$(m,q,d)$-low individual degree test}
is stated in \Cref{fig:test}.
\end{definition}

{
\floatstyle{boxed} 
\restylefloat{figure}
\begin{figure}
With probability~$\tfrac{1}{3}$ each, perform one of the following three tests.
\begin{enumerate}
	\item \textbf{Axis-parallel lines test:}
		Pick a uniformly random role $\br \sim \{\mathrm{A},\mathrm{B}\}$.
		Let $\bu \sim \F_q^m$ be a uniformly random point.
		Select $\bi \sim \{1, \ldots, m\}$ uniformly at random,
		and let $\bell = \{\bu + t \cdot e_{\bi} \mid t \in \F_q\}$
		be the axis-parallel line which passes through~$\bu$ in the $\bi$-th direction.
		\begin{itemize}
		\item[$\circ$] Player~$\br$: 
			Give $\bell$;
			receive the univariate degree-$d$ polynomial $\boldf:\bell \rightarrow \F_q$.
		\item[$\circ$] Player~$\overline{\br}$:
			Give $\bu$;
			receive $\ba \in \F_q$.
		\end{itemize}
		Accept if $\boldf(\bu) = \ba$.
	\item \textbf{Self-consistency test:}
		Let $\bu \sim \F_q^m$ be a uniformly random point.
		\begin{itemize}
		\item[$\circ$] Player~$\mathrm{A}$: 
			Give $\bu$;
			receive $\ba \in \F_q$.
		\item[$\circ$] Player~$\mathrm{B}$:
			Give $\bu$;
			receive $\bb \in \F_q$.
		\end{itemize}
		Accept if $\ba = \bb$.
	\item \textbf{Diagonal lines test:}
		Pick a uniformly random role $\br \sim \{\mathrm{A},\mathrm{B}\}$.
		Let $\bu \sim \F_q^m$ be a uniformly random point.
		Select $\bi \sim \{1, \ldots, m\}$ uniformly at random,
		and let $\bv \in \F_q^m$ be a uniformly random point whose last $m - \bi$ coordinates are~$0$.
		Finally,
		let $\bell = \{\bu + t \cdot \bv \mid t \in \F_q\}$
		be the line which passes through~$\bu$ in direction~$\bv$.
		\begin{itemize}
		\item[$\circ$] Player~$\br$: 
			Give $\bell$;
			receive the univariate degree-$md$ polynomial $\boldf:\bell \rightarrow \F_q$.
		\item[$\circ$] Player~$\overline{\br}$:
			Give $\bu$;
			receive $\ba \in \F_q$.
		\end{itemize}
		Accept if $\boldf(\bu) = \ba$.
	\end{enumerate}
	\caption{The $(m,q,d)$-low individual degree test.\label{fig:test}}
\end{figure}
}

An important class of strategies are those which are \emph{symmetric}.
In this case, the bipartite state $\ket{\psi}$ Alice and Bob share is symmetric.
Furthermore, for any question Alice and Bob receive, they apply the same measurement to their share of the state.
This means that rather than, for example, keeping track of a separate points measurement for Alice and Bob,
we can use a single measurement to refer to both of their strategies,
and similarly for the axis-parallel lines and diagonal lines measurements.
This is formalized in the following definition, where we also consider strategies which are projective.

\begin{definition}[Symmetric, projective strategy]
A \emph{symmetric, projective strategy} for the $(m,q, d)$-low individual degree test
is a tuple $(\psi, A, B, L)$ defined as follows.
\begin{itemize}
\item[$\circ$] $\ket{\psi} \in \calH \ot \calH$ is a bipartite, permutation-invariant state.
\item[$\circ$] $A = \{A^u_a\}$ contains a matrix for each $u \in \F_q^m$ and $a \in \F_q$.
			For each $u$, $A^u$ is a projective measurement on~$\calH$.
\item[$\circ$] $B = \{B^{\ell}_f\}$ contains a matrix for each axis-parallel line~$\ell$ in $\F_q^m$
			and univariate degree-$d$ polynomial $f:\ell\rightarrow \F_q$.
			For each $\ell$, $B^{\ell}$ is a projective measurement on~$\calH$.
\item[$\circ$] $L = \{L^{\ell}_f\}$ contains a matrix for each line~$\ell$ in $\F_q^m$
			and univariate degree-$md$ polynomial $f:\ell\rightarrow \F_q$.
			For each $\ell$, $L^{\ell}$ is a projective measurement on~$\calH$.
\end{itemize}
\end{definition}

We can also consider more general strategies which are no longer assumed to be symmetric.
In this case, Players~$\mathrm{A}$ and~$\mathrm{B}$ each have their own versions of the measurements~$A$, $B$, and~$L$.

\begin{definition}[General projective strategy]
A \emph{projective strategy} for the $(m,q, d)$-low individual degree test
is a tuple $(\psi, A^{\mathrm{A}},B^{\mathrm{A}},L^{\mathrm{A}},A^{\mathrm{B}},B^{\mathrm{B}},L^{\mathrm{B}})$ defined as follows.
\begin{itemize}
\item[$\circ$] $\ket{\psi} \in \calH_{\mathrm{A}} \ot \calH_{\mathrm{B}}$ is a bipartite state.
\end{itemize}
Furthermore, for each~$w \in \{\mathrm{A}, \mathrm{B}\}$:
\begin{itemize}
\item[$\circ$] $A^{w} = \{A^{w,u}_a\}$ contains a matrix for each $u \in \F_q^m$ and $a \in \F_q$.
			For each $u$, $A^{w,u}$ is a projective measurement on~$\calH_{w}$.
\item[$\circ$] $B^{w} = \{B^{w,\ell}_f\}$ contains a matrix for each axis-parallel line~$\ell$ in $\F_q^m$
			and univariate degree-$d$ polynomial $f:\ell\rightarrow \F_q$.
			For each $\ell$, $B^{w,\ell}$ is a projective measurement on~$\calH_{w}$.
\item[$\circ$] $L^w = \{L^{w,\ell}_f\}$ contains a matrix for each line~$\ell$ in $\F_q^m$
			and univariate degree-$md$ polynomial $f:\ell\rightarrow \F_q$.
			For each $\ell$, $L^{w,\ell}$ is a projective measurement on~$\calH_{w}$.
\end{itemize}
\end{definition}

\begin{remark}
Throughout this work, we will only consider projective strategies.
So we will henceforth use the terms ``strategy" and ``symmetric strategy'' to refer exclusively to projective strategies and symmetric, projective strategies, respectively.

In addition, we will spend the vast majority of this work dealing solely with symmetric strategies, as they are notationally simpler to work with.
Much of this work will focus on proving~\Cref{thm:main-induction},
a variant of our main theorem for symmetric strategies.
Our main theorem for general strategies, \Cref{thm:main-formal} below, will be proven in \Cref{sec:induction} by a standard reduction to the symmetric case.
We will only see these more cumbersome-to-write general strategies
in this section, where we state \Cref{thm:main-formal},
and in \Cref{sec:induction}, where we carry out the reduction.
\end{remark}

\begin{definition}[Good strategy]
A strategy is $(\eps, \delta, \gamma)$-good if it
passes the axis-parallel lines test with probability at least $1-\eps$,
the self consistency test with probability at least~$1-\delta$,
and the diagonal lines test with probability at least~$1-\gamma$.
\end{definition}

\begin{remark}\label{rem:good-strat-characterization}
Using notation which will be introduced in \Cref{sec:comparing-measurements} below,
a symmetric strategy is $(\eps, \delta, \gamma)$-good
if and only if it satisfies the following three conditions.
For $\bell$ and $\bu$ as in the axis-parallel lines test,
\begin{equation*}
A^u_a \ot I \simeq_{\eps} I \ot B^{\ell}_{[f(u)=a]},
\quad
\text{and}
\quad
A^u_a \ot I \simeq_{\delta} I \ot A^u_a.
\end{equation*}
And for $\bell$ and $\bu$ as in the diagonal lines test,
\begin{equation*}
A^u_a \ot I \simeq_{\gamma} I \ot L^{\ell}_{[f(u)=a]}.
\end{equation*}
\end{remark}

\begin{notation}\label{not:conditioned-on-last-direction}
Our proof will be via induction,
i.e.\ proving soundness of the $(m+1,q,d)$-low individual degree test
using the soundness of the $(m,q,d)$-low individual degree test.
To do this, we will frequently use the axis-parallel line test in the specific case of $i = m+1$.
Thus, it will be convenient to introduce the following notation.
Let $(\psi, A, B, L)$ be a symmetric strategy for the $(m+1, q, d)$-low individual degree test,
Then for each $u \in \F_q^m$ we will write $B^u_f$ as shorthand for $B^{\ell}_f$,
where $\ell = \{(u, x) \mid x \in \F_q\}$.
For a function $f:\ell \rightarrow \F_q$,
we will also sometimes write $f(x)$ as shorthand for $f(u,x)$.
\end{notation}

\begin{definition}
Consider the $(m,q,d)$-low individual degree test.
For $j \in \{1, \ldots, m\}$,
we refer to the \emph{$j$-restricted diagonal lines test}
as the diagonal lines test conditioned on $\bi = j$.
For example, in the $m$-restricted diagonal lines test,
the line $\bell$ is simply a uniformly random line in $\F_q^m$.
\end{definition}

Now we state our main theorem using notation which will be introduced in \Cref{sec:prelims} below.
This the formal version of \Cref{thm:main-informal}.

\begin{theorem}[Main theorem; quantum soundness of the low individual degree test]\label{thm:main-formal}
  Consider a projective strategy $(\psi, A^{\mathrm{A}},B^{\mathrm{A}},L^{\mathrm{A}},A^{\mathrm{B}},B^{\mathrm{B}},L^{\mathrm{B}})$ 
   which passes the  $(m,q,d)$-low individual degree test with probability at least $1-\eps$.
  Let $k \geq md$ be an integer.
  Let
  \begin{equation*}
  \nu = 100000 k^2m^4 \cdot \Big(\eps^{1/40000} + (d/q)^{1/40000} + e^{-k/(2560000m^2)}\Big).
  \end{equation*}
Then there exists projective measurements $G^{\mathrm{A}}, G^{\mathrm{B}} \in \polymeas{m}{q}{d}$ with the following properties:
\begin{enumerate}
  \item (Consistency with~$A$):  On average over $\bu \sim \F_q^{m}$,
	  \begin{align*}
	  A^{\mathrm{A},u}_a \otimes I
	  	&\simeq_{\nu} I \otimes G^{\mathrm{B}}_{[g(u)=a]},\\
	I \otimes A^{\mathrm{B},u}_a
	  	&\simeq_{\nu}  G^{\mathrm{A}}_{[g(u)=a]} \ot I.
	  \end{align*}
\item (Self-consistency):
	\begin{equation*}
	G^{\mathrm{A}}_g \ot I \simeq_{\nu} I \ot G^{\mathrm{B}}_g.
	\end{equation*}
\end{enumerate}
\end{theorem}

We note that there is a tradeoff in \Cref{thm:main-formal} specified by the parameter~$k$.
As~$k$ increases, $\nu$'s prefactor $k^2$ increases.
On the other hand, as $k$ increases, the term $e^{- k/(2560000m^2)}$ which also occurs in $\nu$ \emph{decreases}.
Thus, when applying this theorem, one must select~$k$ to balance these competing demands.
Typically, choosing $k = \poly(m)$ should be more than sufficient for applications.
We believe that this parameter~$k$ is an artifact of our proof,
and we hope that it will be removed in the future.

%%% Local Variables:
%%% mode: latex
%%% TeX-master: "multilinearity"
%%% End:
